\documentclass[aspectratio=43]{beamer}

% Text packages to stop warnings
\usepackage{lmodern}
\usepackage{textcomp}
\usepackage{ulem}
\usepackage[utf8]{inputenc}
\usepackage[T1]{fontenc}

\usepackage{listings}
\usepackage{tikz}
\usetikzlibrary{arrows,decorations.pathreplacing,positioning}

% Themes
\usetheme{Boadilla}
\setbeamertemplate{footline}[page number]{}
\setbeamertemplate{navigation symbols}{}

% Suppress the navigation bar
\beamertemplatenavigationsymbolsempty

\newenvironment{changemargin}[1]{% 
  \begin{list}{}{% 
    \setlength{\topsep}{0pt}% 
    \setlength{\leftmargin}{#1}% 
    \setlength{\rightmargin}{1em}
    \setlength{\listparindent}{\parindent}% 
    \setlength{\itemindent}{\parindent}% 
    \setlength{\parsep}{\parskip}% 
  }% 
  \item[]}{\end{list}} 

\lstset{basicstyle=\scriptsize, frame=single}

\lstdefinelanguage{JavaScript}{
  keywords={typeof, new, true, false, catch, function, return, null, catch, switch, var, if, in, while, do, else, case, break},
  keywordstyle=\color{blue}\bfseries,
  ndkeywords={class, export, boolean, throw, implements, import, this},
  ndkeywordstyle=\color{darkgray}\bfseries,
  identifierstyle=\color{black},
  sensitive=false,
  comment=[l]{//},
  morecomment=[s]{/*}{*/},
  commentstyle=\color{purple}\ttfamily,
  stringstyle=\color{red}\ttfamily,
  morestring=[b]',
  morestring=[b]"
}

\title{Lecture 29---Compiler Optimizations}
\subtitle{ECE 459: Programming for Performance}
\date{March 18, 2014}

\begin{document}
%%%%%%%%%%%%%%%%%%%%%%%%%%%%%%%%%%%%%%%%%%%%%%%%%%%%%%%%%%%%%%%%%%%%%%%%%%%%%%%%
\begin{frame}[plain]
  \titlepage
\end{frame}
%%%%%%%%%%%%%%%%%%%%%%%%%%%%%%%%%%%%%%%%%%%%%%%%%%%%%%%%%%%%%%%%%%%%%%%%%%%%%%%%

%%%%%%%%%%%%%%%%%%%%%%%%%%%%%%%%%%%%%%%%%%%%%%%%%%%%%%%%%%%%%%%%%%%%%%%%%%%%%%%%
\begin{frame}
  \frametitle{Last Time}
  \begin{changemargin}{1cm}
    Investigated impact of using abstractions.\\[1em]
    Today: how compilers actually make code faster.
  \end{changemargin}
\end{frame}
%%%%%%%%%%%%%%%%%%%%%%%%%%%%%%%%%%%%%%%%%%%%%%%%%%%%%%%%%%%%%%%%%%%%%%%%%%%%%%%%

%%%%%%%%%%%%%%%%%%%%%%%%%%%%%%%%%%%%%%%%%%%%%%%%%%%%%%%%%%%%%%%%%%%%%%%%%%%%%%%%
\part{Compiler Optimizations}
\frame{\partpage}

%%%%%%%%%%%%%%%%%%%%%%%%%%%%%%%%%%%%%%%%%%%%%%%%%%%%%%%%%%%%%%%%%%%%%%%%%%%%%%%%
\begin{frame}
  \frametitle{Introduction}

  \begin{changemargin}{2cm}
    We'll be looking at compiler optimizations.\\[1em]
    Mostly related to performance.\\
    Better {\tt gcc} than you:
      \begin{itemize}
        \item For you, probably a waste of time
        \item Plus: optimizations make your code less readable.
      \end{itemize}
~\\

    Compilers have a host of optimization options.\\
     We'll look at {\tt gcc}.
  \end{changemargin}
\end{frame}
%%%%%%%%%%%%%%%%%%%%%%%%%%%%%%%%%%%%%%%%%%%%%%%%%%%%%%%%%%%%%%%%%%%%%%%%%%%%%%%%

%%%%%%%%%%%%%%%%%%%%%%%%%%%%%%%%%%%%%%%%%%%%%%%%%%%%%%%%%%%%%%%%%%%%%%%%%%%%%%%%
\begin{frame}
  \frametitle{Compiler Optimization}

  \begin{changemargin}{2cm}
    ``Optimization'' is a bit of a misnomer:\\
      \qquad compilers don't generate ``optimal'' code.\\[1em]

    Compilers do generate \structure{better} code, though.\\[1em]

    The program you wrote is too slow. \\
    Contract of a compiler: \\
    \qquad produce a program with same behaviour, but faster.
  \end{changemargin}
\end{frame}
%%%%%%%%%%%%%%%%%%%%%%%%%%%%%%%%%%%%%%%%%%%%%%%%%%%%%%%%%%%%%%%%%%%%%%%%%%%%%%%%

%%%%%%%%%%%%%%%%%%%%%%%%%%%%%%%%%%%%%%%%%%%%%%%%%%%%%%%%%%%%%%%%%%%%%%%%%%%%%%%%
\begin{frame}
  \frametitle{gcc Optimization Levels}

  \begin{changemargin}{2cm}
  {\tt -O1} ({\tt -O})

  \begin{itemize}
    \item Reduce code size and execution time.
    \item No optimizations that increase compilation time.
  \end{itemize}

  {\tt -O2}

  \begin{itemize}
    \item All optimizations except space vs. speed tradeoffs.
  \end{itemize}
  
  {\tt -O3}

  \begin{itemize}
    \item All optimizations.
  \end{itemize}

  {\tt -O0} (default)

  \begin{itemize}
    \item Fastest compilation time, debugging works as expected.
  \end{itemize}
  \end{changemargin}
\end{frame}
%%%%%%%%%%%%%%%%%%%%%%%%%%%%%%%%%%%%%%%%%%%%%%%%%%%%%%%%%%%%%%%%%%%%%%%%%%%%%%%%

%%%%%%%%%%%%%%%%%%%%%%%%%%%%%%%%%%%%%%%%%%%%%%%%%%%%%%%%%%%%%%%%%%%%%%%%%%%%%%%%
\begin{frame}[fragile]
  \frametitle{Disregard Standards, Acquire Speedup}

  \begin{changemargin}{2cm}
  {\tt -Ofast}

  \begin{itemize}
    \item All {\tt -O3} optimizations and non-standards compliant optimizations,
      particularly {\tt -ffast-math}. \\ 
 (Like {\tt -fast} on Solaris.)
  \end{itemize}
  This turns off exact implementations of IEEE or ISO rules/specifications for math
  functions.\\[1em]

  Generally, if you don't care about the exact result, you can use this for
  a speedup
  \end{changemargin}
\end{frame}
%%%%%%%%%%%%%%%%%%%%%%%%%%%%%%%%%%%%%%%%%%%%%%%%%%%%%%%%%%%%%%%%%%%%%%%%%%%%%%%%

\section{Scalar Optimizations}
%%%%%%%%%%%%%%%%%%%%%%%%%%%%%%%%%%%%%%%%%%%%%%%%%%%%%%%%%%%%%%%%%%%%%%%%%%%%%%%%
\begin{frame}[fragile]
  \frametitle{Constant Folding}

  \begin{changemargin}{2cm}
  \begin{lstlisting}
i = 1024 * 1024
  \end{lstlisting}

  Why do something later you can do now?\\[1em]

  The compiler will not emit code that does the multiplication at runtime.\\
  It will simply use the computed value

  \begin{lstlisting}
i = 1048576
  \end{lstlisting}

  \begin{itemize}
    \item Enabled at all optimization levels.
  \end{itemize}
  \end{changemargin}
\end{frame}
%%%%%%%%%%%%%%%%%%%%%%%%%%%%%%%%%%%%%%%%%%%%%%%%%%%%%%%%%%%%%%%%%%%%%%%%%%%%%%%%

%%%%%%%%%%%%%%%%%%%%%%%%%%%%%%%%%%%%%%%%%%%%%%%%%%%%%%%%%%%%%%%%%%%%%%%%%%%%%%%%
\begin{frame}[fragile]
  \frametitle{Common Subexpression Elimination}

  \begin{changemargin}{1.5cm}

  {\tt -fgcse}

  \begin{itemize}
    \item Perform a global\footnote{across basic blocks} common subexpression elimination pass.
    \item (Also performs global constant and copy propagation.)
    \item Enabled at {\tt -O2}, {\tt -O3}.
  \end{itemize}
  ~\\[1em]

  {\bf Example:} 

  \begin{lstlisting}
a = b * c + g;
d = b * c * d;
  \end{lstlisting}

  Instead of computing {\tt b * c} twice, we compute it once, and reuse the
  value in each expression
  \end{changemargin}
\end{frame}
%%%%%%%%%%%%%%%%%%%%%%%%%%%%%%%%%%%%%%%%%%%%%%%%%%%%%%%%%%%%%%%%%%%%%%%%%%%%%%%%

%%%%%%%%%%%%%%%%%%%%%%%%%%%%%%%%%%%%%%%%%%%%%%%%%%%%%%%%%%%%%%%%%%%%%%%%%%%%%%%%
\begin{frame}[fragile]
  \frametitle{Constant Propagation}

  \begin{changemargin}{2cm}

  Moves constant values from definition to use.

  \begin{itemize}
    \item Valid if there's no intervening redefinition.
  \end{itemize}
  \vfill
  {\bf Example:}
  \begin{lstlisting}
int x = 14;
int y = 7 - x / 2;
return y * (28 / x + 2);
  \end{lstlisting}

  with constant propagation would produce:

  \begin{lstlisting}
int x = 14;
int y = 0;
return 0;
  \end{lstlisting}
  \end{changemargin}
\end{frame}
%%%%%%%%%%%%%%%%%%%%%%%%%%%%%%%%%%%%%%%%%%%%%%%%%%%%%%%%%%%%%%%%%%%%%%%%%%%%%%%%

%%%%%%%%%%%%%%%%%%%%%%%%%%%%%%%%%%%%%%%%%%%%%%%%%%%%%%%%%%%%%%%%%%%%%%%%%%%%%%%%
\begin{frame}[fragile]
  \frametitle{Copy Propagation}

  \begin{changemargin}{2cm}
  Replaces direct assignments with their values. \\
  Usually runs after
  common subexpression elimination.\\[1em]

  {\bf Example:}
  \begin{lstlisting}
y = x
z = 3 + y
  \end{lstlisting}

  with copy propagation would produce:

  \begin{lstlisting}
z = 3 + x
  \end{lstlisting}
  \end{changemargin}
\end{frame}
%%%%%%%%%%%%%%%%%%%%%%%%%%%%%%%%%%%%%%%%%%%%%%%%%%%%%%%%%%%%%%%%%%%%%%%%%%%%%%%%

%%%%%%%%%%%%%%%%%%%%%%%%%%%%%%%%%%%%%%%%%%%%%%%%%%%%%%%%%%%%%%%%%%%%%%%%%%%%%%%%
\begin{frame}[fragile]
  \frametitle{Dead Code Elimination}

  \begin{changemargin}{2cm}
  {\tt -fdce}

  \begin{itemize}
    \item Remove any code guaranteed not to execute.
    \item Enabled at all optimization levels.
  \end{itemize}
~\\

  {\bf Example:}
  \begin{lstlisting}
if (0) {
    z = 3 + x;
}
  \end{lstlisting}
  would not be included in the final executable
  \end{changemargin}
\end{frame}
%%%%%%%%%%%%%%%%%%%%%%%%%%%%%%%%%%%%%%%%%%%%%%%%%%%%%%%%%%%%%%%%%%%%%%%%%%%%%%%%

\section{Loop Optimizations}
%%%%%%%%%%%%%%%%%%%%%%%%%%%%%%%%%%%%%%%%%%%%%%%%%%%%%%%%%%%%%%%%%%%%%%%%%%%%%%%%
\begin{frame}[fragile]
  \frametitle{Loop Unrolling}

  \begin{changemargin}{2cm}
  {\tt -funroll-loops}
  \begin{itemize}
    \item Unroll loops, using a fixed number of iterations.
  \end{itemize}
  \vfill
  {\bf Example:}

  \begin{lstlisting}
for (int i = 0; i < 4; ++i)
    f(i)
  \end{lstlisting}

would be transformed to:

  \begin{lstlisting}
f(0)
f(1)
f(2)
f(3)
  \end{lstlisting}
(Also useful for SIMD).
  \end{changemargin}
\end{frame}
%%%%%%%%%%%%%%%%%%%%%%%%%%%%%%%%%%%%%%%%%%%%%%%%%%%%%%%%%%%%%%%%%%%%%%%%%%%%%%%%

%%%%%%%%%%%%%%%%%%%%%%%%%%%%%%%%%%%%%%%%%%%%%%%%%%%%%%%%%%%%%%%%%%%%%%%%%%%%%%%%
\begin{frame}[fragile]
  \frametitle{Loop Interchange}

  \begin{changemargin}{2cm}
  {\tt -floop-interchange}\\
  Enhances locality; big wins possible.\\[1em]

  {\bf Example:} in C the following:
  \begin{lstlisting}
for (int i = 0; i < N; ++i)
    for (int j = 0; j < M; ++j)
        a[j][i] = a[j][i] * c
  \end{lstlisting}
  would be transformed to this:
  \begin{lstlisting}
for (int j = 0; j < M; ++j)
    for (int i = 0; i < N; ++i)
        a[j][i] = a[j][i] * c
  \end{lstlisting}
  since C is {\bf row-major} (meaning a[1][1] is beside a[1][2]).\\
  (Other
  possibility is {\bf column-major}.)
  \end{changemargin}
\end{frame}
%%%%%%%%%%%%%%%%%%%%%%%%%%%%%%%%%%%%%%%%%%%%%%%%%%%%%%%%%%%%%%%%%%%%%%%%%%%%%%%%

%%%%%%%%%%%%%%%%%%%%%%%%%%%%%%%%%%%%%%%%%%%%%%%%%%%%%%%%%%%%%%%%%%%%%%%%%%%%%%%%
\begin{frame}[fragile]
  \frametitle{Loop Fusion}

  \begin{changemargin}{2cm}
  {\bf Example:}
  \begin{lstlisting}
for (int i = 0; i < 100; ++i)
    a[i] = 4

for (int i = 0; i < 100; ++i)
    b[i] = 7
  \end{lstlisting}
  would be transformed to this:
  \begin{lstlisting}
for (int i = 0; i < 100; ++i) {
    a[i] = 4
    b[i] = 7
}
  \end{lstlisting}
  \vfill
  Trade-off  between data locality and loop overhead; \\ 
  the inverse is also an optimization, called {\bf loop fission}.
  \end{changemargin}
\end{frame}
%%%%%%%%%%%%%%%%%%%%%%%%%%%%%%%%%%%%%%%%%%%%%%%%%%%%%%%%%%%%%%%%%%%%%%%%%%%%%%%%

%%%%%%%%%%%%%%%%%%%%%%%%%%%%%%%%%%%%%%%%%%%%%%%%%%%%%%%%%%%%%%%%%%%%%%%%%%%%%%%%
\begin{frame}[fragile]
  \frametitle{Loop-Invariant Code Motion}
  
  \begin{changemargin}{2cm}
  \begin{itemize}
    \item Moves invariants out of a loop.
    \item Also called loop hoisting.
  \end{itemize}

  {\bf Example:}
  \begin{lstlisting}
for (int i = 0; i < 100; ++i) {
    s = x * y;
    a[i] = s * i;
}
  \end{lstlisting}
  would be transformed to this:
  \begin{lstlisting}
s = x * y;
for (int i = 0; i < 100; ++i) {
    a[i] = s * i;
}
  \end{lstlisting}
~\\
  Reduces the amount of work we have to do for each iteration of the loop.
  \end{changemargin}
\end{frame}
%%%%%%%%%%%%%%%%%%%%%%%%%%%%%%%%%%%%%%%%%%%%%%%%%%%%%%%%%%%%%%%%%%%%%%%%%%%%%%%%

\section{Other Optimizations}
%%%%%%%%%%%%%%%%%%%%%%%%%%%%%%%%%%%%%%%%%%%%%%%%%%%%%%%%%%%%%%%%%%%%%%%%%%%%%%%%
\begin{frame}
  \frametitle{Devirtualization}

  \begin{changemargin}{2cm}
  {\tt -fdevirtualize}

  \begin{itemize}
    \item Attempt to convert virtual function calls to direct calls.
    \item Enabled with {\tt -O2}, {\tt -O3}.
  \end{itemize}
~\\[1em]

  Virtual functions impose overhead and impede other optimizations.\\[1em]
  C++: must read
  the object's vtable and branch to
  the correct function.
  \end{changemargin}
\end{frame}
%%%%%%%%%%%%%%%%%%%%%%%%%%%%%%%%%%%%%%%%%%%%%%%%%%%%%%%%%%%%%%%%%%%%%%%%%%%%%%%%

%%%%%%%%%%%%%%%%%%%%%%%%%%%%%%%%%%%%%%%%%%%%%%%%%%%%%%%%%%%%%%%%%%%%%%%%%%%%%%%%
\begin{frame}[fragile]
  \frametitle{Devirtualization (example)}
  \begin{changemargin}{2cm}

  {\bf Example:}

  \begin{lstlisting}
class A {
    virtual void m();
};

class B : public A {
    virtual void m();
}

int main(int argc, char *argv[]) {

    std::unique_ptr<A> t(new B);
    t.m();
}
  \end{lstlisting}
  Devirtualization could eliminate RTTI access; \\ \qquad instead just call to B's {\tt m}.\\
  Needs call graph, which can be tricky to compute (more soon).
  \end{changemargin}
\end{frame}
%%%%%%%%%%%%%%%%%%%%%%%%%%%%%%%%%%%%%%%%%%%%%%%%%%%%%%%%%%%%%%%%%%%%%%%%%%%%%%%%

%%%%%%%%%%%%%%%%%%%%%%%%%%%%%%%%%%%%%%%%%%%%%%%%%%%%%%%%%%%%%%%%%%%%%%%%%%%%%%%%
\begin{frame}[fragile]
  \frametitle{Scalar Replacement of Aggregates}

\begin{center}
  \includegraphics[width=0.8\textwidth]{L18/Censored_rubber_stamp.png}
\end{center}
\end{frame}
%%%%%%%%%%%%%%%%%%%%%%%%%%%%%%%%%%%%%%%%%%%%%%%%%%%%%%%%%%%%%%%%%%%%%%%%%%%%%%%%

%%%%%%%%%%%%%%%%%%%%%%%%%%%%%%%%%%%%%%%%%%%%%%%%%%%%%%%%%%%%%%%%%%%%%%%%%%%%%%%%
\begin{frame}
  \frametitle{Aliasing and Pointer Analysis}

  \begin{changemargin}{1.5cm}

     We've seen {\tt restrict}: tell compiler that variables don't
      alias.\\[1em]

    {\it Pointer analysis} automatically tracks the variables in your program to
      determine whether or not they alias.\\[1em]
    
    If they don't alias, we can reason about side effects, reorder
    accesses, and do other types of optimizations.
  \end{changemargin}

\end{frame}
%%%%%%%%%%%%%%%%%%%%%%%%%%%%%%%%%%%%%%%%%%%%%%%%%%%%%%%%%%%%%%%%%%%%%%%%%%%%%%%%

%%%%%%%%%%%%%%%%%%%%%%%%%%%%%%%%%%%%%%%%%%%%%%%%%%%%%%%%%%%%%%%%%%%%%%%%%%%%%%%%
\begin{frame}
  \frametitle{Call Graph}
  
  \begin{changemargin}{1.5cm}
\begin{itemize}
\item     A directed graph showing relationships between functions.
\end{itemize}
~\\
     Relatively simple to compute in C\\\qquad (function pointers complicate things).\\

     Hard for virtual function calls (C++/Java).\\[1em]

     Virtual calls require pointer analysis to disambiguate.
  \end{changemargin}
\end{frame}
%%%%%%%%%%%%%%%%%%%%%%%%%%%%%%%%%%%%%%%%%%%%%%%%%%%%%%%%%%%%%%%%%%%%%%%%%%%%%%%%

%%%%%%%%%%%%%%%%%%%%%%%%%%%%%%%%%%%%%%%%%%%%%%%%%%%%%%%%%%%%%%%%%%%%%%%%%%%%%%%%
\begin{frame}[fragile]
  \frametitle{Importance of Call Graphs}

  \begin{changemargin}{2cm}

  The call graph can enable optimization of the following:

  \begin{lstlisting}
int n;

int f() { /* opaque */ }

int main() {
    n = 5;
    f();
    printf("%d\n", n);
}
  \end{lstlisting}
~\\
  We could propagate the constant value 5 in {\tt main()},\\ \qquad as long as {\tt f()}
  does not write to {\tt n}.
  \end{changemargin}
\end{frame}
%%%%%%%%%%%%%%%%%%%%%%%%%%%%%%%%%%%%%%%%%%%%%%%%%%%%%%%%%%%%%%%%%%%%%%%%%%%%%%%%

%%%%%%%%%%%%%%%%%%%%%%%%%%%%%%%%%%%%%%%%%%%%%%%%%%%%%%%%%%%%%%%%%%%%%%%%%%%%%%%%
\begin{frame}[fragile]
  \frametitle{Tail Recursion Elimination}

  \begin{changemargin}{2cm}

  {\tt -foptimize-sibling-calls}

  \begin{itemize}
    \item Optimize sibling and tail recursive calls.
    \item Enabled at {\tt -O2} and {\tt -O3}.
  \end{itemize}

  {\bf Example:}

  \begin{lstlisting}
int bar(int N) {
    if (A(N))
        return B(N);
    else
        return bar(N);
}
  \end{lstlisting}

  Compiler can just replace the call to {\tt bar} by a {\tt goto}.\\
  Avoids function call overhead 
  and reduces call stack use.
  \end{changemargin}
\end{frame}
%%%%%%%%%%%%%%%%%%%%%%%%%%%%%%%%%%%%%%%%%%%%%%%%%%%%%%%%%%%%%%%%%%%%%%%%%%%%%%%%

%%%%%%%%%%%%%%%%%%%%%%%%%%%%%%%%%%%%%%%%%%%%%%%%%%%%%%%%%%%%%%%%%%%%%%%%%%%%%%%%
\begin{frame}[fragile]
  \frametitle{Branch Predictions}

  \begin{changemargin}{2cm}
     gcc attempts to guess the probability of each branch to 
      best order the code.\\
     (for an {\tt if}, fall-through is most efficient, why?)\\[1em]

     Use {\tt \_\_builtin\_expect(expr, value)} to help GCC, if
      you know the run-time characteristics of your program.\\[1em]
  {\bf Example (in the Linux kernel):}

  \begin{lstlisting}
#define likely(x)       __builtin_expect((x),1)
#define unlikely(x)     __builtin_expect((x),0)
  \end{lstlisting}
  \end{changemargin}
\end{frame}
%%%%%%%%%%%%%%%%%%%%%%%%%%%%%%%%%%%%%%%%%%%%%%%%%%%%%%%%%%%%%%%%%%%%%%%%%%%%%%%%

%%%%%%%%%%%%%%%%%%%%%%%%%%%%%%%%%%%%%%%%%%%%%%%%%%%%%%%%%%%%%%%%%%%%%%%%%%%%%%%%
\begin{frame}
  \frametitle{Architecture-Specific}

  \begin{changemargin}{1.5cm}
  Two common ones: {\tt -march} and {\tt -mtune}\\ \qquad\qquad ({\tt -march} implies {\tt -mtune}).

  \begin{itemize}
    \item enable  specific instructions that not all CPUs  support\\
      \qquad \qquad (SSE4.2, etc.)
    \vfill
    \item {\bf Example:} {\tt -march=corei7}
    \vfill
    \item Good to use on your local machine, not so much for shipped code.
  \end{itemize}
  \end{changemargin}
\end{frame}
%%%%%%%%%%%%%%%%%%%%%%%%%%%%%%%%%%%%%%%%%%%%%%%%%%%%%%%%%%%%%%%%%%%%%%%%%%%%%%%%

\part{Profile-Guided Optimization}
\frame{\partpage}

%%%%%%%%%%%%%%%%%%%%%%%%%%%%%%%%%%%%%%%%%%%%%%%%%%%%%%%%%%%%%%%%%%%%%%%%%%%%%%%%
\begin{frame}
  \frametitle{Moving Beyond Purely-Static Approaches}

  \begin{changemargin}{2cm}
    So far: purely-static (compile-time) optimizations.\\[1em]

    Static approach has limitations: 
    \begin{itemize}
      \item how aggressively to inline?
      \item which branch to predict?
    \end{itemize}
    \structure{Dynamic} information answers these questions better.

  \end{changemargin}

\end{frame}
%%%%%%%%%%%%%%%%%%%%%%%%%%%%%%%%%%%%%%%%%%%%%%%%%%%%%%%%%%%%%%%%%%%%%%%%%%%%%%%%

%%%%%%%%%%%%%%%%%%%%%%%%%%%%%%%%%%%%%%%%%%%%%%%%%%%%%%%%%%%%%%%%%%%%%%%%%%%%%%%%
\begin{frame}
  \frametitle{Profile-Guided Optimization Workflow}

  \begin{changemargin}{2cm}
\begin{enumerate}
\item \structure{Compile} the program; \\
\hspace*{1em} produce an instrumented binary.
\item \structure{Run} the instrumented binary:\\
\hspace*{1em} on a representative test suite;\\
\hspace*{1em} get run-time profiling information.\\
\item \structure{Compile} final version, with the profiles you've
collected.
\end{enumerate}

  \end{changemargin}

\end{frame}
%%%%%%%%%%%%%%%%%%%%%%%%%%%%%%%%%%%%%%%%%%%%%%%%%%%%%%%%%%%%%%%%%%%%%%%%%%%%%%%%

%%%%%%%%%%%%%%%%%%%%%%%%%%%%%%%%%%%%%%%%%%%%%%%%%%%%%%%%%%%%%%%%%%%%%%%%%%%%%%%%
\begin{frame}
  \frametitle{Price of Performance from Profile-Guided Optimization}
  \begin{changemargin}{2cm}
\begin{tabbing}
    ~\!(--) \hspace*{.4em} \= Complicates the compilation process.\\
    (+) \> Produces faster code (if inputs are good).\\[1em]
\end{tabbing}

    How much faster? 5--25\%, per Microsoft.\\[1em]

 Solaris, Microsoft VC++, and GNU gcc all
 support profile-guided optimization (to some extent).

  \end{changemargin}

\end{frame}
%%%%%%%%%%%%%%%%%%%%%%%%%%%%%%%%%%%%%%%%%%%%%%%%%%%%%%%%%%%%%%%%%%%%%%%%%%%%%%%%

%%%%%%%%%%%%%%%%%%%%%%%%%%%%%%%%%%%%%%%%%%%%%%%%%%%%%%%%%%%%%%%%%%%%%%%%%%%%%%%%
\begin{frame}
  \frametitle{Another Option}
  \begin{changemargin}{2cm}
    Just-in-time compilers (like Java Virtual Machines)
    automatically do profile-guided optimizations:
\begin{itemize}
\item compile code on-the-fly anyway;
\item always collecting measurements for the current run.
\end{itemize}


  \end{changemargin}

\end{frame}
%%%%%%%%%%%%%%%%%%%%%%%%%%%%%%%%%%%%%%%%%%%%%%%%%%%%%%%%%%%%%%%%%%%%%%%%%%%%%%%%

%%%%%%%%%%%%%%%%%%%%%%%%%%%%%%%%%%%%%%%%%%%%%%%%%%%%%%%%%%%%%%%%%%%%%%%%%%%%%%%%
\begin{frame}
  \frametitle{What's Profiling Data Good For?}
  \begin{changemargin}{2cm}
Synergistic optimizations:
\begin{itemize}
\item \structure{Inlining}: inlining rarely-used code is a net lose; \\
inlining often-executed code is a win, \\
\qquad enables more optimizations.\\[1em]
Everyone does this.\\[1em]

\item \structure{Improving Cache Locality}: make commonly-called functions, basic blocks adjacent;\\
improves branch prediction (below) and with {\tt switch}.

\end{itemize}


  \end{changemargin}

\end{frame}
%%%%%%%%%%%%%%%%%%%%%%%%%%%%%%%%%%%%%%%%%%%%%%%%%%%%%%%%%%%%%%%%%%%%%%%%%%%%%%%%

%%%%%%%%%%%%%%%%%%%%%%%%%%%%%%%%%%%%%%%%%%%%%%%%%%%%%%%%%%%%%%%%%%%%%%%%%%%%%%%%

\begin{frame}
  \frametitle{More Uses of Profiling Data}
  \begin{changemargin}{1cm}

\begin{itemize}
\item \structure{Branch Prediction}: a common architecture assumption---\\
\qquad forward branches not taken, backwards branches taken.\\
can make these assumptions more true with profiles;\\
for a forward branch, make sure common case is fall-through.\\[1em]

\item \structure{Virtual Call Prediction}: inline the most common target \\ \qquad (guarded by a conditional).
\end{itemize}
  \end{changemargin}

\end{frame}
%%%%%%%%%%%%%%%%%%%%%%%%%%%%%%%%%%%%%%%%%%%%%%%%%%%%%%%%%%%%%%%%%%%%%%%%%%%%%%%%

\part{Summary}
\frame{\partpage}

%%%%%%%%%%%%%%%%%%%%%%%%%%%%%%%%%%%%%%%%%%%%%%%%%%%%%%%%%%%%%%%%%%%%%%%%%%%%%%%%
\begin{frame}
  \frametitle{Summary}

  \begin{changemargin}{1cm}
    Got a feel for what the optimization levels do.\\
    We saw examples of compiler optimizations. For instance:
\begin{itemize}
\item Scalar optimizations. \\
(e.g. common subexpression elimination,\\ \quad constant~propagation, copy propagation).
\item Redundant code optimizations.\\
(e.g. dead code elimination).
\item Loop optimizations.\\
(e.g. loop interchange, loop fusion, loop fission, \\ \quad software~pipelining).
\item Alias and pointer analysis; call graphs.\\
(uses: devirtualization, inlining, tail recursion elimination).
\end{itemize}
  \end{changemargin}


\end{frame}
%%%%%%%%%%%%%%%%%%%%%%%%%%%%%%%%%%%%%%%%%%%%%%%%%%%%%%%%%%%%%%%%%%%%%%%%%%%%%%%%

%%%%%%%%%%%%%%%%%%%%%%%%%%%%%%%%%%%%%%%%%%%%%%%%%%%%%%%%%%%%%%%%%%%%%%%%%%%%%%%%
\begin{frame}
  \frametitle{Summary II}

  \begin{changemargin}{0.5cm}
     Full list of {\tt gcc} options: \\ ~~~{\small \url{http://gcc.gnu.org/onlinedocs/gcc/Optimize-Options.html}}\\
      More optimization examples: \\ ~~\url{http://www.digitalmars.com/ctg/ctgOptimizer.html}\\[1em]
  \end{changemargin}
  \begin{changemargin}{1.5cm}

     Also, profile-guided optimization:\\
\qquad see what actually happens at runtime, optimize for that.
  \end{changemargin}
\end{frame}
%%%%%%%%%%%%%%%%%%%%%%%%%%%%%%%%%%%%%%%%%%%%%%%%%%%%%%%%%%%%%%%%%%%%%%%%%%%%%%%%

\end{document}

